\section{Systemübersicht} 
\subsection{Produktperspektive}
Das Nothaltesystem ist ein sicherheitskritisches Subsystem des autonomen Miniatur-LKWs.
Es arbeitet parallel zur Fahrzeugsteuerung und überwacht kontinuierlich Sensorwerte sowie Operatorbefehle.
Bei Erreichen definierter Gefährdungszustände greift es mit höchster Priorität ein und überschreibt alle anderen Fahrfunktionen.
Das System interagiert über ROS, CAN-Kommunikation und Aktor-Schnittstellen mit bestehenden Komponenten wie Lokalisierung, Fahrsteuerung, LiDAR-Sensorik und Motorregelung.
Es ersetzt nicht die autonome Fahrlogik, sondern ergänzt diese um eine deterministische Notfall- und Sicherheitsmechanik.

\subsection{Systemkontext}
\label{sec:systemkontext}

Der Systemkontext beschreibt alle externen Komponenten, Akteure und Subsysteme, die mit dem Nothaltesystem interagieren.

\subsubsection{Akteur}

\begin{table}[H]
  \centering
  \caption{Menschlicher Akteur des Nothaltesystems}
  \label{tab:akteur_nothalt}
  \begin{tabularx}{\textwidth}{lX}
    \toprule
    \textbf{Akteur} & \textbf{Beschreibung} \\
    \midrule
    Fahrer/Operator &
    Der Operator kann über verschiedene Eingabegeräte Steuerbefehle,
    Softstop-Signale und Fahrkommandos an das Nothaltesystem bzw.
    die Fahrzeugsteuerung senden. Er ist der primäre
    Mensch–Maschine-Interaktionspartner. \\
    \bottomrule
  \end{tabularx}
\end{table}

\subsubsection{Externe Systemkomponenten}

\begin{table}[h]
  \centering
  \caption{Externe Systemkomponenten im Kontext des Nothaltesystems}
  \label{tab:systemkomponenten_nothalt}
  \begin{tabularx}{\textwidth}{lX}
    \toprule
    \textbf{Komponente} & \textbf{Rolle / Interaktion} \\
    \midrule
    Fernbedienung &
    Sendet Fahrbefehle und Softstop-Signale an das Nothaltesystem oder das autonome Fahrsystem. \\

    Autonomes Fahrsystem &
    Übermittelt Fahrbefehle und erhält den Nothaltstatus zurück. \\

    \texttt{ArduinoMaster} &
    Übergibt den softwareseitigen Nothaltstatus an das hardwarenahe System. \\

    \texttt{ArduinoSlave} &
    Setzt Stop-Signale in PWM-Signale um und steuert die Motoren entsprechend. \\

    Antrieb / Motorcontroller &
    Führt Brems- bzw. Stoppbefehle physikalisch aus. \\

    UI / RViz &
    Visualisiert Nothaltstatus, Warnungen und ermöglicht ggf. Softstop-Befehle. \\

    Logging / ROS-Bags &
    Zeichnet Nothalt-Ereignisse und Statusmeldungen für Analysezwecke auf. \\

    Dachwarnleuchte &
    Erhält ein An-/Aus-Signal zur optischen Warnung bei aktivem Nothalt. \\

    Sensorik (LiDAR, Encoder) &
    Liefert Umgebungsdaten zur autonomen Erkennung von Gefahrensituationen. \\
    \bottomrule
  \end{tabularx}
\end{table}

\subsection{Schnittstellenübersicht}
\label{sec:schnittstellenuebersicht}

Der Systemkontext weist mehrere klar abgegrenzte Schnittstellen zwischen Nothaltesystem, Fahrzeugsteuerung, Sensorik und Peripherie auf. Die wichtigsten logischen Schnittstellen sind in Tabelle~\ref{tab:schnittstellen_nothalt} zusammengefasst.

\begin{table}[h]
  \centering
  \caption{Hauptschnittstellen des Nothaltesystems}
  \label{tab:schnittstellen_nothalt}
  \begin{tabular}{p{0.25\textwidth}p{0.65\textwidth}}
    \hline
    \textbf{Schnittstelle} & \textbf{Beschreibung} \\
    \hline
    Fahrbefehle &
    Von Fernbedienung oder autonomem Fahrsystem in Richtung
    Nothaltesystem. Die Fahrbefehle umfassen u.\,a.\ Geschwindigkeits-,
    Richtungs- und Betriebsmodusvorgaben. \\

    Nothaltstatus (ROS $\rightarrow$ Arduino) &
    Übergibt sicherheitskritische Nothaltinformationen vom ROS-basierten
    Nothaltesystem an das Arduino-System, das anschließend den
    physischen Stopp auslöst (z.\,B.\ Abschalten der Motorfreigabe). \\

    CAN-Signale &
    Weiterleitung von E-Stop-Informationen zu den relevanten
    Fahrzeugsteuerungen über den CAN-Bus, damit auch andere
    Steuergeräte den Nothaltzustand berücksichtigen können. \\

    Warnsignal (Dachleuchte) &
    Einfaches Ausgangssignal zur optischen Kennzeichnung eines aktiven
    Nothalts. Bei gesetztem Nothaltstatus wird die Dachwarnleuchte
    dauerhaft aktiviert. \\

    Sensorströme (Lidar, Encoder) &
    Eingehende Sensordatenströme, die zur automatisierten
    Hinderniserkennung sowie zur Bewertung der Fahrzeugbewegung
    verwendet werden. Dazu zählen insbesondere 2D-/3D-Lidar-Scans
    und Encoder-Rückmeldungen. \\

    Logging-Daten &
    Ereignisdaten und Sensordaten, die zur späteren Analyse,
    Nachvollziehbarkeit und Validierung des Sicherheitsverhaltens
    persistiert werden (z.\,B.\ als ROS-Bag-Dateien). \\
    \hline
  \end{tabular}
\end{table}

Durch diese Abgrenzung wird deutlich, dass das Nothaltesystem als
zentrale sicherheitskritische Instanz zwischen Bediener, Fahrzeugregelung,
Sensorik und sicherheitsrelevanten Aktoren fungiert.


